\section{The soft acoustic limit}\label{app:D.37}

For a long-wavelength acoustic phonon the phonon-assisted amplitude \cref{eq:13.2-M} tends to the limit \cref{eq:13.6-soft}, of order $\ell_{\vect{Q}\nu} \propto |\vect{Q}|^{-1/2}$, although its intraband terms are of the form $0/0$ and taken alone would suggest $\ell_{\vect{Q}\nu}/|\vect{Q}|$.

\subsection*{A Ward identity for the amplitude}

Let $\op{\mathfrak{t}}_{\vect{Q}}$ be any one-electron operator that raises the crystal momentum by $\vect{Q}$, with matrix elements $\mathfrak{t}_{mn}(\vect{\kappa},\vect{Q}) = \langle m,\vect{\kappa}+\vect{Q}|\op{\mathfrak{t}}_{\vect{Q}}|n,\vect{\kappa}\rangle$, and define the vertex
\begin{equation}
  \delta g_{mn}(\vect{\kappa},\vect{Q}) = \bigl(\varepsilon_{m,\vect{\kappa}+\vect{Q}} - \varepsilon_{n\vect{\kappa}} \pm \omega_{\vect{Q}\nu}\bigr)\,\mathfrak{t}_{mn}(\vect{\kappa},\vect{Q}) ,
  \label{eq:D.37-dg}
\end{equation}
with the sign of emission (upper) or absorption (lower) as in \cref{eq:13.2-M}. It is the matrix element of the commutator of $\op{\mathfrak{t}}_{\vect{Q}}$, multiplied by the phonon operator, with the Hamiltonian of free electrons and phonons. Insert $\delta g$ in place of $g$ into \cref{eq:13.2-M}. In the first term, with $p$ at $\vect{k}+\vect{q}$, the vertex is $\delta g_{cp}(\vect{k}+\vect{q},\vect{Q}) = (\varepsilon_{c\vect{k}'} - \varepsilon_{p,\vect{k}+\vect{q}} \pm \omega_{\vect{Q}\nu})\mathfrak{t}_{cp}$, and on the energy shell, $\omega = \varepsilon_{c\vect{k}'} - \varepsilon_{v\vect{k}} \pm \omega_{\vect{Q}\nu}$, the denominator $\omega - \varepsilon_{p,\vect{k}+\vect{q}} + \varepsilon_{v\vect{k}}$ is the same number. Their ratio is exactly 1. In the second term the vertex is $\delta g_{pv}(\vect{k},\vect{Q}) = (\varepsilon_{p,\vect{k}+\vect{Q}} - \varepsilon_{v\vect{k}} \pm \omega_{\vect{Q}\nu})\mathfrak{t}_{pv}$ and the denominator its negative, so the ratio is $-1$. Both statements hold band by band, including the intraband terms $p = c$ in the first sum and $p = v$ in the second, whose denominators are of order $|\vect{Q}|$: the vertex vanishes with them, and no quotient $0/0$ ever has to be taken. The sums over all bands $p$ are then complete, and
\begin{equation}
  M^\pm[\delta g] = \sum_p\bigl[\mathfrak{t}_{cp}(\vect{k}+\vect{q},\vect{Q})\,\rho_{pv}(\vect{k},\vect{q}) - \rho_{cp}(\vect{k}+\vect{Q},\vect{q})\,\mathfrak{t}_{pv}(\vect{k},\vect{Q})\bigr]
  = \langle c,\vect{k}'|\,[\op{\mathfrak{t}}_{\vect{Q}},\ee^{\ii\vect{q}\cdot\vect{r}}]\,|v,\vect{k}\rangle .
  \label{eq:D.37-contact}
\end{equation}
Since $M^\pm$ is linear in the vertex, for any $\op{\mathfrak{t}}_{\vect{Q}}$
\begin{equation}
  M^\pm[g] = M^\pm[g - \delta g] + \langle c,\vect{k}'|\,[\op{\mathfrak{t}}_{\vect{Q}},\ee^{\ii\vect{q}\cdot\vect{r}}]\,|v,\vect{k}\rangle .
  \label{eq:D.37-ward}
\end{equation}
This is exact on the energy shell, at every $\vect{Q}$, for every band and branch. It is a Ward identity in the sense of \cref{sec:8.3}: a transformation of the states generated by $\op{\mathfrak{t}}_{\vect{Q}}$ changes the electron--phonon vertex by a commutator with the unperturbed Hamiltonian and the density operator by a commutator with $\op{\mathfrak{t}}_{\vect{Q}}$, and the physical amplitude does not change. Off the shell the ratios are no longer $\pm1$ and the identity fails, as the band-sum theorem of \cref{app:D.35} does.

\subsection*{The comoving vertex}

A long acoustic wave displaces the lattice near $\vect{r}$ by $\vect{u}(\vect{r}) = (\ell_{\vect{Q}\nu}/\sqrt N)\,\hat{\vect{e}}_\nu\,\ee^{\ii\vect{Q}\cdot\vect{r}}B_{\vect{Q}\nu}$ (\cref{sec:6.2}), and an orbital carried along with it changes by $-\vect{u}\cdot\nabla\psi = -\ii\,\vect{u}\cdot\op{\vect{p}}\,\psi$. Take $\op{\mathfrak{t}}_{\vect{Q}}$ to be minus this change per unit of $N^{-1/2}B_{\vect{Q}\nu}$, symmetrized so that it is defined for a $\vect{u}$ that varies in space,
\begin{equation}
  \op{\mathfrak{t}}_{\vect{Q}} = \ii\,\ell_{\vect{Q}\nu}\,\tfrac12\bigl\{\ee^{\ii\vect{Q}\cdot\vect{r}},\,\hat{\vect{e}}_\nu\cdot\op{\vect{p}}\bigr\} .
  \label{eq:D.37-s}
\end{equation}
With a nonlocal pseudopotential, $\op{\vect{p}}$ is the generator of translations of \cref{sec:6.2}, which moves the nonlocal projectors with the ions.

\emph{Contact term.} Since $[\op{\vect{p}},\ee^{\ii\vect{q}\cdot\vect{r}}] = \vect{q}\,\ee^{\ii\vect{q}\cdot\vect{r}}$, the commutator in \cref{eq:D.37-ward} is $\ii\ell_{\vect{Q}\nu}(\hat{\vect{e}}_\nu\cdot\vect{q})\,\ee^{\ii(\vect{q}+\vect{Q})\cdot\vect{r}}$, and
\begin{equation}
  \langle c,\vect{k}'|\,[\op{\mathfrak{t}}_{\vect{Q}},\ee^{\ii\vect{q}\cdot\vect{r}}]\,|v,\vect{k}\rangle = \ii\,\ell_{\vect{Q}\nu}\,(\hat{\vect{e}}_\nu\cdot\vect{q})\,\rho_{cv}(\vect{k},\vect{q}+\vect{Q}) .
  \label{eq:D.37-recoil}
\end{equation}
It is the term of first order in the displacement of $\ee^{\ii\vect{q}\cdot(\vect{r} - \vect{u})}$, the density of electrons that ride on the displaced lattice. The commutator of $\op{\vect{p}}$ with $\ee^{\ii\vect{q}\cdot\vect{r}}$, which is $\vect{q}\,\ee^{\ii\vect{q}\cdot\vect{r}}$, is the electron part of the translation algebra behind \cref{eq:8.3-mass}: the density at wavevector $\vect{q}$ carries momentum $\vect{q}$.

\emph{No rigid translation.} At $\vect{Q} = 0$, since $\omega_{\vect{Q}\nu} \to 0$ on an acoustic branch,
\begin{equation*}
  \begin{gathered}
  \mathfrak{t}_{mn}(\vect{\kappa},0) = \ii\ell_{\vect{0}\nu}\,\hat{\vect{e}}_\nu\cdot\vect{p}_{mn}(\vect{\kappa}) , \\
  \delta g_{mn}(\vect{\kappa},0) = \ii\ell_{\vect{0}\nu}\,(\varepsilon_{m\vect{\kappa}} - \varepsilon_{n\vect{\kappa}})\,\hat{\vect{e}}_\nu\cdot\vect{p}_{mn}(\vect{\kappa}) .
  \end{gathered}
\end{equation*}
For the acoustic polarization $(M_\kappa/M_{\mathrm{cell}})^{1/2}\hat{\vect{e}}_\nu$, the vertex \cref{eq:A.6-g} at $\vect{Q} = 0$ is $\ell\,\hat{\vect{e}}_\nu\cdot\sum_\kappa\langle m|\partial_{\kappa,\vect{0}}V_{\mathrm{KS}}|n\rangle$, which \cref{eq:6.2-asr} evaluates to the same expression. The comoving vertex
\begin{equation}
  g^{\mathrm{cm}}_{mn\nu}(\vect{\kappa},\vect{Q}) = g_{mn\nu}(\vect{\kappa},\vect{Q}) - \delta g_{mn}(\vect{\kappa},\vect{Q})
  \label{eq:D.37-gcm}
\end{equation}
therefore vanishes at $\vect{Q} = 0$ for every pair of bands, interband as well as intraband. The lab-frame vertex does not: its interband elements at $\vect{Q} \to 0$ are those of a rigid translation, of order $\ell_{\vect{Q}\nu} \propto |\vect{Q}|^{-1/2}$.

\emph{First order in $\vect{Q}$.} Write \cref{eq:D.37-gcm} as $[g - (\varepsilon_{m,\vect{\kappa}+\vect{Q}} - \varepsilon_{n\vect{\kappa}})\mathfrak{t}_{mn}] \mp \omega_{\vect{Q}\nu}\mathfrak{t}_{mn}$. The bracket is the potential seen by an electron in coordinates that move with the lattice. Over a region short compared with $2\pi/|\vect{Q}|$ the wave is a uniform strain there, $u_{\alpha\beta} = \frac{\ii}{2}(Q_\alpha\hat e_{\nu\beta} + Q_\beta\hat e_{\nu\alpha})\ell_{\vect{Q}\nu}$ per unit amplitude, and in comoving coordinates the reduced wavevector of a Bloch state does not change. The intraband element of the bracket is therefore the band shift under a uniform strain at fixed reduced wavevector, $\sum D^n_{\alpha\beta}u_{\alpha\beta}$, which is the Bardeen--Shockley argument of \cref{sec:6.2} and defines the deformation-potential tensor $D^n$ of \cref{eq:6.2-D}. The last term is $\mp c_\nu|\vect{Q}|$ times $\mathfrak{t}_{mn} \to \ii\ell\,\hat{\vect{e}}_\nu\cdot\vect{p}_{mn}$: the coupling to the velocity of the moving lattice. To first order in $\vect{Q}$, with $\hat{\vect{Q}} = \vect{Q}/|\vect{Q}|$,
\begin{equation}
  g^{\mathrm{cm}}_{nn\nu}(\vect{\kappa},\vect{Q}) = \ii\,\ell_{\vect{Q}\nu}|\vect{Q}|\bigl[\hat{\vect{Q}}\cdot D^n\cdot\hat{\vect{e}}_\nu \mp c_\nu\,\hat{\vect{e}}_\nu\cdot\vect{p}_{nn}(\vect{\kappa})\bigr] ,
  \label{eq:D.37-gcmnn}
\end{equation}
and the interband elements of $g^{\mathrm{cm}}$ are also of order $\ell_{\vect{Q}\nu}|\vect{Q}|$. Adding $\delta g_{nn} = (|\vect{Q}|\,\hat{\vect{Q}}\cdot\vect{v}^{\mathrm{g}}_n \pm c_\nu|\vect{Q}|)\,\ii\ell\,\hat{\vect{e}}_\nu\cdot\vect{p}_{nn}$ back, the velocity terms cancel, and the intraband vertex in the laboratory frame is
\begin{equation}
  g_{nn\nu}(\vect{\kappa},\vect{Q}) = \ii\,\ell_{\vect{Q}\nu}|\vect{Q}|\;\hat{\vect{Q}}\cdot\bigl(D^n + \vect{v}^{\mathrm{g}}_n\otimes\vect{p}_{nn}\bigr)\cdot\hat{\vect{e}}_\nu ,
  \label{eq:D.37-Xi}
\end{equation}
with $\vect{v}^{\mathrm{g}}_n$ the group velocity at $\vect{\kappa}$. At a band extremum $\vect{v}^{\mathrm{g}}_n = 0$, and \cref{eq:D.37-Xi} reduces to \cref{eq:6.2-D}. Away from it the laboratory vertex carries the Galilean term $\vect{v}^{\mathrm{g}}_n\otimes\vect{p}_{nn}$, equal to $\vect{p}_{nn}\otimes\vect{p}_{nn}/m_e$ for a local potential, which the frozen-strain calculation of \cref{app:E.17} does not produce, because it computes $D^n$ alone; a nonlocal potential, with $\vect{v}^{\mathrm{g}}_n \ne \vect{p}_{nn}/m_e$, adds a part that is not symmetric in the two indices.

\subsection*{The limit}

Apply \cref{eq:D.37-ward} with the comoving choice. The contact term is \cref{eq:D.37-recoil}, with $\rho_{cv}(\vect{k},\vect{q}+\vect{Q}) \to \rho_{cv}(\vect{k},\vect{q})$ as $\vect{Q}\to0$. In $M^\pm[g^{\mathrm{cm}}]$, the interband terms have vertices of order $\ell|\vect{Q}|$ and denominators that stay finite, and they are of order $\ell|\vect{Q}|$. Two terms are not: the intraband terms, $p = c$ in the first sum and $p = v$ in the second, whose denominators vanish with $\vect{Q}$. On the shell they are
\begin{equation*}
  \omega - \varepsilon_{c,\vect{k}+\vect{q}} + \varepsilon_{v\vect{k}} = \varepsilon_{c\vect{k}'} - \varepsilon_{c,\vect{k}+\vect{q}} \pm \omega_{\vect{Q}\nu} = |\vect{Q}|\bigl(\hat{\vect{Q}}\cdot\vect{v}^{\mathrm{g}}_c \pm c_\nu\bigr) + O(Q^2) ,
\end{equation*}
with $\vect{v}^{\mathrm{g}}_c$ at $\vect{k}+\vect{q}$, and $\varepsilon_{v\vect{k}} - \varepsilon_{v,\vect{k}+\vect{Q}} \mp \omega_{\vect{Q}\nu} = -|\vect{Q}|(\hat{\vect{Q}}\cdot\vect{v}^{\mathrm{g}}_v \pm c_\nu) + O(Q^2)$, with $\vect{v}^{\mathrm{g}}_v$ at $\vect{k}$. Divided into \cref{eq:D.37-gcmnn}, the factor $|\vect{Q}|$ cancels and each term stays of order $\ell_{\vect{Q}\nu}$. Collecting the contact term and the two intraband terms gives \cref{eq:13.6-soft}, with a remainder $O(\ell_{\vect{Q}\nu}|\vect{Q}|)$. The leading term is of order $|\vect{Q}|^{-1/2}$ and the correction of order $|\vect{Q}|^{1/2}$, a relative correction of order $|\vect{Q}|$. The expansion of the intraband detunings to first order requires $|\vect{Q}|$ small compared with the scale on which the group velocity changes, $|\hat{\vect{Q}}\cdot\vect{v}^{\mathrm{g}}_n \pm c_\nu|$ divided by the band curvature, which at an extremum is $|\vect{Q}| \ll 2m^*c_\nu$ for a band of mass $m^*$. For a degenerate band, such as the valence maximum of silicon, the intraband terms are matrices in the degenerate partners, and the $O(Q)$ splitting along $\hat{\vect{Q}}$ is diagonalized first. \Cref{we:13-soft} checks each step in a crystal of rigid ions in one dimension, where the laboratory vertex is computed directly from the displaced ions.

The same result follows from \cref{eq:13.2-M} in the laboratory frame, without \cref{eq:D.37-ward}. There the rigid-translation part of the interband vertices, $\ii\ell(\varepsilon_m - \varepsilon_n)\hat{\vect{e}}_\nu\cdot\vect{p}_{mn}$, combines over all bands, by the completeness used in \cref{eq:D.37-contact}, into $\ii\ell[\hat{\vect{e}}_\nu\cdot(\vect{q} - \vect{p}_{cc} + \vect{p}_{vv})]\rho_{cv}$, and the intraband terms with the vertex \cref{eq:D.37-Xi} supply the rest; collecting the $\vect{p}$ terms gives \cref{eq:13.6-soft} again. Setting the interband vertices of an acoustic mode to zero near $\vect{Q} = 0$ therefore loses exactly that combination, the recoil term and the bare Galilean terms.

\subsection*{Inside and outside the window}

On the energy shell, $\vect{Q}$ cannot be small inside the off-shell window. A final state with $\omega \le E_{\mathrm{min}}(\vect{q}) - \delta_{\mathrm{w}}$ and an electron at $\vect{k}' = \vect{k}+\vect{q}+\vect{Q}$ has $\varepsilon_{c,\vect{k}+\vect{q}} - \varepsilon_{v\vect{k}} \ge E_{\mathrm{min}}(\vect{q})$, so $\varepsilon_{c,\vect{k}+\vect{q}} - \varepsilon_{c\vect{k}'} \ge \delta_{\mathrm{w}} \pm \omega_{\vect{Q}\nu}$. Emission gives $|\vect{Q}| \ge \delta_{\mathrm{w}}/(v^{\mathrm{g}}_{\max} - c_\nu)$ and absorption the weaker bound, which holds for either sign,
\begin{equation}
  |\vect{Q}| \ge \frac{\delta_{\mathrm{w}}}{v^{\mathrm{g}}_{\max} + c_\nu} ,
  \label{eq:D.37-Qmin}
\end{equation}
with $v^{\mathrm{g}}_{\max}$ the largest conduction-band group velocity. Outside the window $\vect{Q}$ can go to zero, and the emission of soft phonons by the real pair is integrable (\cref{sec:13.6}), with $n_{\vect{Q}\nu} \simeq T/c_\nu|\vect{Q}|$ at $T > 0$.
